\begin{abstract}

    3D-printed auxetic metamaterials have gained significant interest in recent years due to the ease of fast prototyping. While previous studies examine the nonlinear properties of the auxetic metamaterials from the geometrical perspective, including unit cell shape, direction, type, and magnitude, none of the studies closely examine how the nonlinear properties change with 3D printing parameters, including infill patterns, layer height, layer thickness, infill density, and so on. Hence, this study aims to closely examine how the nonlinear properties of the auxetic metamaterials vary with 3D printing parameters and to develop a mathematical model to predict these relationships using artificial neural networks and to inversely design the auxetic metamaterials.

\end{abstract}